\cotefile{64}
\scanpage{15}
\mspage{7}

Considérons le groupe $\mathcal D$ en tant que groupe ``abstrait'', il est \unsure{canoniqt|canoniquement} extension de $\mathcal D_{\mathrm{ab}}=W$ ($\simeq\mathbb F_2\times\mathbb F_2$ non canoniqt) par $\mathbb F_2$. Dans $\mathcal D_{\mathrm{ab}}$ il y a un élément non nul distingué $\overline\omega$, \deleted{\unsure{image inverse|image réciproque}} \inserted{dont l'image} inverse est l'ens $\underline\omega_{\unsure{Q|\mathrm{le}}}$ des deux éléments d'ordre $4$ de $\mathcal D$. On voit aisément que les deux autres éléments \unsure{deux|non} nuls de $\mathcal D_{\mathrm{ab}}$ sont interchangés par le groupe des auto. de $\mathcal D$, on trouve donc
\diagram{\begin{tikzcd}[ampersand replacement=\&,column sep=small]
1\arrow[r,dash] \& \check W\arrow[r] \& \operatorname{Aut}\mathcal D\arrow[r] \& \pm1\arrow[r,dash] \& 1\\[-2pt]
\& {\substack{\text{autom de}\\\deleted{\unsure{\mathcal D}}\text{l'ext. }\mathcal D\text{ de}\\ W\text{ par }\mathbb F_2\text{ qui}\\\text{induisent l'identité}\\\text{dans }W\text{ et }\mathbb F_2}} \& \& {\underbrace{\vphantom{\pm1}\hspace{1em}}_{\substack{\text{interchange}\\\text{les les deux}\\\text{éléments de}\\\deleted{\unsure{\mathcal D_{\mathrm{ab}}}}W\setminus\{0,\overline\omega\}}}} \&
\end{tikzcd}}

% slanted
\margin{\uline{NB} on a en fait un iso. can. $\check W\simeq W$ (valable pour tt vectoriel de rg $2$ sur $\mathbb F_2$)}

En fait $\operatorname{Aut}\mathcal D$ est produit $\frac12$ direct de $\pm1$ par $\check W$, $\pm1$ opérant par échange des \inserted{$2$} éléments non \deleted{\unsure{distins|distin}} privilégiés de $\check W\simeq W$. Quand on tient compte que $\mathcal D$ provient d'un carré, donc il y a choix d'un élément $\underline c$ de $W\setminus\{0,\overline\omega\}$ (correspondant \deleted{\unsure{au}} à $C\deleted{\unsure{=|\simeq}}V_S\setminus\{1,\underline a\}$), l'autre élément étant $\underline d$ (correspondant à $D=V_A\setminus\{1,\underline a\}$). Les autom de $\mathcal D$ qui respectent cette structure \deleted{forment} le sous-groupe $\check W\simeq W$ de $\operatorname{Aut}\mathcal D$, formé des autom \inserted{d'ext.} qui induisent l'identité sur le quotient $W$. Ce sont aussi les autom intérieurs, qui forment bien le groupe $W\simeq\mathcal D/\underbrace{\{1,a\}}_{\text{centre}}$.

Comme l'hom $\mathcal D\xrightarrow{\mathrm{int}}\operatorname{Aut}\mathcal D$ n'est pas injectif, on voit que la donnée du \deleted{\unsure{carré|autre}} groupe $\mathcal D=\operatorname{Aut}Q$ (qui revient à la donnée de $C$ et de $D$, \inserted{cf plus bas}) \deleted{\unsure{comme on a} $\mathcal D=\mathbb F_2^C\mathop{\times}_{\mathbb F_2}\mathbb F_2^D$)}

% archive sheet number: 14

\scanpage{16}

ne peut suffire à réconstruire $Q$ à isom. unique près -- mais seulement $Q$ ``modulo antipodisme''. Considérant le carré épinglé $Q_0$ et le groupe de ses automorphismes
\[
\deleted{\mathcal D_4}\ \mathbb D_4=\operatorname{Aut}(Q_0)
\]
avec la suite exacte canonique
\diagram{\begin{tikzcd}[ampersand replacement=\&,column sep=small]
1\arrow[r]\&\mathbb F_2\arrow[r]\&\mathbb D_4\arrow[r]\&\mathbb F_2^2\arrow[r,dash]\&1
\end{tikzcd}}
et la donnée de notre $\mathcal D$ revient à celle d'un torseur sous \inserted{le quotient} $\mathbb F_2^2$ \inserted{de $\mathbb D_4$} i.e. de deux ens. à deux éléments (ce sont en fait $D$ et $C$, dans cet ordre \ldots), tandis que celle d'un carré revient à la donnée d'un torseur sur $\mathbb D_4$ lui-$=$.

On a une opération \inserted{(par autom. int.)} de $\mathcal D^+$ sur $D\inserted{\hookrightarrow\mathcal D^+}$, la donnée de $Q$ (correspondant au groupe d'autom $\mathcal D$) \deleted{\unsure{ne}} revient à celle d'un relèvement de \deleted{\unsure{cette}} \deleted{\inserted{\unsure{cet}}} l'\deleted{\unsure{é}}opération \deleted{\unsure{induite de}} en \deleted{\unsure{$\mathcal D^+$|des|les}} une opération \deleted{\unsure{sur}} transitive sur un ens. $S$ \inserted{(donc un esp. hom sous $\mathcal D$)}, avec $S/\underline a\simeq D$. Description duale via $C\hookrightarrow\mathcal D^+\ldots$

\uline{Remarque.} Le groupe $\operatorname{Aut}(\mathbb D_4)$ est isom. à $\mathbb D_4$ (mais \uline{pas} via autom. intérieurs, bien sûr) ce qui implique qu'il y a \inserted{quand même} une équivalence de groupoïdes $(\text{Carrés})\xrightarrow{\approx}(\text{groupes isom. à }\mathbb D_4)$ (mais qui n'est \unsure{pas|$\vDash$} donnée par \deleted{\unsure{$\mathbb D_4$)}} $Q\mapsto\operatorname{Aut}(Q)$ !), qui mériterait d'être comprise \ldots

% archive sheet number: 15

\scanpage{17}
\mspage{8}

Si $\mathcal D$ est un groupe isomorphe à $\mathbb D_4$, l'ens $S$ des éléments d'ordre $2$ \deleted{\unsure{non centraux}} \inserted{$\ne\underline a$ (où} $\underline a\longrightarrow$ l'unique élément central $\ne1$) \deleted{est} a quatre éléments, \deleted{qui forment deux classes par translation} \inserted{sur lequel on a une involution} \inserted{sans} point fixe par
\[
x\longmapsto x+\underline a\overset{\mathrm{def}}=x'
\]
donc à $\mathcal D$ correspond un carré $Q_{\mathcal D}$, dont l'ens des sommets est $S$, et l'\deleted{\unsure{automorphisme de|automorphisme des}} \inserted{action de} l'antipodisme sur $S$ est l'involution précédente. Mais le foncteur de groupoïdes $\mathcal D\mathrel{\text{---}}Q_{\mathcal D}$
\diagram{\begin{tikzcd}[ampersand replacement=\&]
(\text{groupes isom. à }\mathbb D_4)\arrow[r]\&(\text{carrés combinatoires})
\end{tikzcd}}
est une équivalence de groupoïdes, car on reconstitue à partir de $(S,\underline a_S:\inserted{x\mapsto x'})$ $\mathcal D$ par \deleted{\unsure{purement}} la formule
\[
\mathcal D=\bigl\{\underbrace{x\in S,\ \underline a}_{\text{générateurs}}\ \big|\ 
\underbrace{[x,y]}_{\substack{\text{commutateur}\\xyx^{-1}y^{-1}}}=\underline a\ \text{ si }\ y\ne x,x'\bigr\}
\]

% slanted
\margin{\uline{NB} Si à un carré \inserted{$Q$} on associe le groupe $\mathcal D=\operatorname{Aut}Q$, puis le carré associé à ce groupe, on trouve le carré dont l'ens des sommets est $C\amalg D$, l'action de l'antipodisme dessus \deleted{\unsure{étant}} \unsure{comme orbites} $C$ et $D$, i.e. les diagonales sont $C$ et $D$. On trouve donc ``un carré avec diagonale \deleted{\unsure{privil}} distinguée\ldots''}

Si on part, pour décrire la structure de carré, de l'ens des sommets, $S$, cette structure est définie par la donnée d'un ens. \inserted{$S$} à $4$ éléments -- i.e. d'un torseur sous un vectoriel $V_S$ de rang $2$ sur $\mathbb F_2$ -- et d'une \deleted{\unsure{tr}} involution ss pts fixes de $S$, i.e. d'un $\underline a_S\in V_S^*$ \deleted{Mais se} \deleted{\unsure{$V=V_S^*$|$V_S=V_S^*$}} \inserted{(i.e. une} partition de type $(2,2)$ de $S$). Mais pour

% archive sheet number: 16

\scanpage{18}

un vectoriel \inserted{$V$} de rg $2$ sur $\mathbb F_2$, la donnée d'un $a\in\deleted{\unsure{\mathbb F_2|\mathbb F_2^*}}V^*$ revient à la donnée d'une partie de $V$ qui soit une base (savoir $V^*\setminus\{a\}$) ; donc la donnée d'un \inserted{couple $(V,\underline a)$ d'un} vectoriel $V$ de rg $2$ sur $\mathbb F_2$, et d'un élément \inserted{$a$} de $V^*$, revient à la donnée d'un ens. $C$ à deux éléments, par $V\simeq\mathbb F_2^C$, $a$ l'élément diagonal de $\mathbb F_2^C$. Ainsi, la donnée d'un carré \deleted{\unsure{rev}}$Q$ \deleted{\unsure{via les ens.}} (via l'ens des sommets) équivaut à celle d'un couple $(C,S)$, où $C$ \deleted{\unsure{$(\mathbb F_2)$}} est un ens à deux éléments, et $S$ est un torseur sous $\mathbb F_2^C$. Mais la donnée d'un tel torseur revient à celle d'une famille de \inserted{$\mathbb F_2$-}torseurs indexée par $C$, i.e. d'une application de degré $2$
\[
A\longrightarrow C,
\]
% slanted
\margin{ici les \unsure{deux|dits|\deleted{deux}} ens $A_c$ $(c\in C)$ \unsure{apparaissent|se présentent} \deleted{\unsure{comme}} comme \uline{quotients} de $S$, par les relations d'équivalence définies resp. par \deleted{les} \deleted{deux} \inserted{l'élément} \deleted{\unsure{$c$}} \inserted{$c'=c+a$} de $V^*$ (considérés comme involutions ss pt fixes de $S$) Mais je préfère ici prendre la convention inverse, en \deleted{identifiant} intervertissant $A_c$ et $A_{c'}$.}

$A$ a quatre éléments, et est muni d'une involution sans pts fixes (qui donne le passage au quotient vers $C$), par la formule
\[
S\simeq\prod_{c\in C}\deleted{\unsure{\mathbb F_2}} A_c=\text{Ens des sections de } A\text{ sur } C
\]
qui donne en même temps un plongement
\[
S\hookrightarrow\mathscr P_2(A)\unsure{.}
\]
d'où la relation d'incidence entre $S$ et $A$\ldots

% archive sheet number: 17

\scanpage{19}
\mspage{9}

Dualement, \deleted{\unsure{$A$ esp de|$A$ app de}} la donnée d'un carré-arêtes $A$ s'interprète comme la donnée d'un torseur sous un groupe $\mathbb F_2^D$, ce qui équivaut à la donnée d'une famille de \inserted{$\mathbb F_2$-}torseurs $(S_d)_{d\in D}$, \deleted{\unsure{indexés sur $S$}} \inserted{par}
\[
A\simeq\prod_{d\in D}S_d\qquad(\simeq\text{sections de }S\text{ sur }D)
\]
-- la donnée de cette famille équivalent à la donnée d'un ensemble (à quatre éléments) $S$ et d'un morphisme de degré $2$
\[
S\mathrel{\unsure{\text{---}|\longrightarrow}}D
\]
qui avec la formule précédente donne
\[
A\hookrightarrow\mathscr P_2(S)
\]
et la relation d'incidence entre $S$ et $A$\ldots

\uline{Formulaire} On a, pour un carré $Q$, les trois ensembles à deux éléments
\[
\begin{array}{ccc}
 C,&D,&\underline\omega\\
 \text{codiagonales}&\text{diagonales}&\text{orientations}
\end{array}
\]
plus les deux paquets de deux ens. à deux éléments (donc $4$ autres ens. à deux éléments)
\[
\begin{aligned}
S\xrightarrow{\mathrm{deg.}\,2}D,\qquad & A\xrightarrow{\mathrm{deg.}\,2}C\\
(D\simeq S/\underline a_S).\qquad &(C\simeq A/\underline a_A)
\end{aligned}
\]
On a le formulaire d'iso can.

% archive sheet number: 18

\scanpage{20}

\[
\boxed{\begin{aligned}
(1)\quad &\underline\omega_S\simeq C,\qquad \underline\omega_A\simeq D\\
(2)\quad &\underline\omega\wedge C\wedge D\simeq\mathord{1\mkern-4mu\rule{.055em}{.66em}}_{\mathbb F_2}\quad\text{i.e.}\quad
\left\{\begin{aligned}
\underline\omega&\simeq C\wedge D\\ C&\simeq D\wedge\underline\omega\\ D&\simeq\underline\omega\wedge C
\end{aligned}\right.\\
(3)\quad &\deleted{C}D\simeq S/\underline a_S,\qquad C\simeq A/\underline a_A\\
(4)\quad &S\simeq\prod_{c\in C}A_c,\qquad A\simeq\prod_{d\in D} S_d\\
(5)\quad &D\simeq\bigwedge_{c\in C}A_c,\qquad C\simeq\bigwedge_{d\in D}S_d\\
&\text{ces \inserted{deux} dernières relations provenant de}\\
&\underline\omega\simeq\left(\bigwedge_{c\in C}A_c\right)\wedge C\quad
\text{d'où}\quad\underbrace{\underline\omega\wedge C}_{D\text{ par }(2)}
\simeq\bigwedge_{c\in C}A_c.
\end{aligned}}
\]

3. Complément sur les \unsure{$2$|$3$}-ss-groupes \deleted{de Sylow} d'un groupe $\inserted{G=}\unsure{\mathfrak S|G}_I$ $(\operatorname{card} I=4)$.

% slanted
\margin{À commencer par les $3$-sous-groupes de Sylow}

Soit $V$ le sous-groupe distingué $\simeq\mathbb F_2^2$, de sorte que $V^*\simeq\mathscr P_{2,2}(I)$. On a
{\scriptsize
\diagram{\begin{tikzcd}[ampersand replacement=\&,column sep=small,row sep=6em]
{\substack{\text{Ens des structures}\\\text{de Carré sur }Q\\
(\text{et }Q\text{ correspond à }x\\\text{qui correspond à }\mathcal D,\text{ où}\\
\operatorname{Aut}Q=\mathcal D)}}\simeq V^*
\arrow[rr,leftrightarrow,"{\substack{\text{pour tt ss}\\\text{gr. cyclique }\mathcal D^+\\\sim}}","{\substack{\text{il a un}\\\text{unique él. d'}\\\text{ordre 2, qui est}\\\text{dans }V^*\text{, et}\\\mathcal D^+\text{ en est le centralisateur}}}"' ]\arrow[dr,leftrightarrow,"\unsure{2|\sim}"']
\& \hspace{17ex}
\& {\substack{\text{Sous-groupes cycliques }\mathcal D^+\\\text{d'ordre 4 de }G}}\arrow[dl,leftrightarrow,"\sim"]\\
{\substack{\text{pour tt sous-groupe}\\\text{de Sylow }\mathcal D,\text{ il y a}\\
\text{un unique élément}\\ x\text{ central }\ne1,\text{ et celui}\\
\text{ci est dans }V^*\text{, et }\mathcal D\\\text{est le normalisateur}\\
\text{ou centralisateur de }x}}
\& {\substack{\unsure{3|2}\text{-sous-groupes}\\\text{de Sylow de }G}}
\& {\substack{\text{tout sous-groupe de}\\\text{Sylow }\mathcal D\text{ a un unique ss-}\\
\text{groupe cyclique d'ordre 4,}\\\mathcal D^+\text{, et il en est le}\\\text{normalisateur}}}
\end{tikzcd}}}

% archive sheet number: 19

\scanpage{21}
\mspage{10}

Quels sont les sous-groupes de $\inserted{G=}\unsure{\mathfrak S|G}_I$ \deleted{\unsure{$e$}}, $\ne e$ et $\mathfrak S_I$ ?

A) Sous-groupes d'ordre une puissance de $2$, i.e. contenus dans un $2$-sous-groupe de Sylow $\mathcal D$.

\deleted{Les sous-groupes cycliques correspondent aux}

1) Ss-groupes d'ordre $2$, correspondent aux éléments d'ordre $2$ :

\margin{$V^*\qquad3$}a) les \inserted{trois} éléments de $V^*$, qui forment une orbite \unsure{sous|\deleted{sous}} $\operatorname{Int}G$

\margin{$\widetilde V^*\qquad6$}b) les six transpositions ou réflexions, qui forment une autre orbite sous $\unsure{\operatorname{Int}|\deleted{\operatorname{Int}}}G$, et se groupent en trois paquets (indexés par $V^*$) de deux, par la relation $\tau\tau'\in V^*$, ou par celle d'être contenues dans un $=$ \inserted{$2$-}ss-groupe de Sylow.

% slanted
\margin{i.e. $\tau,\tau'$ ont $=$ image ds $\unsure{\mathfrak S_3|G_3}$}

2) Sous-groupes \deleted{cycliq} d'ordre $4$

\margin{$V^*\qquad3$}c) \inserted{les trois} sous-groupes cycliques d'ordre $4$, qui \deleted{\unsure{se répartissent en un en un}} \inserted{correspondent $1$-$1$ aux trois} sous-groupes de Sylow i.e. les trois structures de carré sur $I$.

\margin{$1$}d) Sous-groupes non cycliques d'ordre $4$, donc isom. à $\mathbb F_2\times\mathbb F_2$ -- il y a d'abord $V^*$

\margin{$V^*\qquad3$}e) puis ceux qui ne sont pas distingués, chacun contenu dans un unique $2$-sous groupe de Sylow, dans lequel il y en \deleted{\unsure{deux}} \inserted{un} comme ça \inserted{(en plus de $V$, qu'on a exclu)}, donc il y en a en tout \deleted{\uline{six}} \inserted{trois} sous-groupes \deleted{$\mathbb F_2\times\mathbb F_2$} $(2,2)$ non distingués, qui \deleted{forment \uuline{deux}} \inserted{forment une} \unsure{classe|classes} de conjugaison \deleted{distinctes, de trois chacune savoir les $V_S\unsure{\subset|\in}\mathcal D$ et les $V_A$, correspondant}

3) Sous-groupes d'ordre $8$ -- ce sont les

\margin{$V^*\qquad3$}f) trois sous-groupes de Sylow de $G$

% archive sheet number: 20

\scanpage{22}

\uline{B)} Considérons ensuite les sous-groupes \inserted{de $G$} d'ordre multiple de $3$, donc d'ordre $3$, $6$ ou $12$.

\margin{$I\qquad4$}g) Sous-groupes d'ordre $3$, ce sont les $4$ $3$-sous-groupes de Sylow, indexés par les éléments de $I$

\margin{$I\qquad4$}h) Les sous-groupes d'ordre $6$ ne sont pas cycliques (il n'y a pas d'élément d'ordre $6$) donc sont isomorphes à $\mathbb D_3$, \deleted{\unsure{ou|on}} ce sont les quatre stabilisateurs des pts de $I$, i.e. les normalisateurs des $3$-sous groupes de Sylow

\margin{$1$}i) Il y a un seul sous-groupe d'ordre $12$ i.e. d'indice $2$, du seul fait que \deleted{\unsure{$\mathfrak S$|$G$}} $G_{\mathrm{ab}}\simeq\{\pm1\}$ et c'est $G^+\simeq\unsure{\mathfrak S|G}_I^+$, ss-groupe \uline{alterné} de $\mathfrak S_I$

Il y a donc \inserted{exactement} $19$ sous-groupes qui sont des deux-sous-groupes, $\ne e$, se répartissant en \uline{six} classes de conjugaison \deleted{\unsure{(i.e.|(aux}} \inserted{(correspondant à} $6$ classes d'\unsure{isomorphismes|isomorphisme} d'espaces homogènes sous $G$), et $9$ sous-groupes \inserted{$\ne G$} qui sont d'ordre multiple de $3$, se répartissant en trois autres classes de conjugaison. Cela fait en tout $28$ sous-groupes, donc $30$ en comptant $\{\unsure{e|1}\}$ et $G$, se répartissant en $9$ classes de conjugaison (resp. $11$ en comp\deleted{tant} $\{e\}$ et $G$), qui correspondent à autant de classes d'isomorphie d'espaces homogènes sous $G$.

\margin{Pour bien faire, il faudrait les \unsure{inclure|rentrer} ds un diagramme d'inclusions, ou diagr. de morphismes d'espaces homogènes.}

% archive sheet number: 21

\scanpage{23}
\mspage{11}

4. \uline{Relations avec le cube orienté}

Si $I$ est un ens. à $4$ éléments, $J\simeq\mathscr P_{2,2}(I)\simeq V_I^*$ l'ens à trois éléments associé, on considère l'application de degré $2$

\diagram{\begin{tikzcd}[ampersand replacement=\&,column sep=small,row sep=0pt]
(*)\&\operatorname{Circ}(I)\arrow[r]\&\unsure{J|\mathbb J}.\\
\&\substack{\text{ens des}\\\text{permutations}\\\text{circulaires}\\\text{sur }I}\&u\mapsto u^2
\end{tikzcd}}

[ \inserted{NB} On a aussi l'application de degré $2$

\diagram{\begin{tikzcd}[ampersand replacement=\&,column sep=small,row sep=0pt]
\mathscr P_2(I)\arrow[r]\&J\\
\substack{\simeq\text{ens. des}\\\text{``transpositions}\\\text{de }I\\\text{(ou ``réflexions'')}}\&
\substack{A\mapsto\{A,I\smallsetminus A\}\\\text{ou }\forall\,\tau\ \text{réflexion,}\\\exists!\,\tau'\ \text{réflexion telle}\\\text{que }\tau\tau'\in V^*}
\end{tikzcd}}

ce sont vraiment des \unsure{revêt\deleted{ement}s|revêt\deleted{ts}s} de nature distincte, les espaces homogènes \inserted{sous $G$} $\operatorname{Circ}(I)$ et $\mathscr P_2(I)$ n'étant pas $\widetilde{\mathrm m}$ isomorphes (le stabilisateur d'un élément de $\operatorname{Circ}$ est $\simeq\mathbb Z/4\mathbb Z$, celui d'un élément de $\mathscr P_2(I)$ est $\simeq\mathbb F_2\times\mathbb F_2$) -- c'est le premier \deleted{des} deux, \inserted{\deleted{\unsure{$\mathscr P$|$\mathfrak P$}} $\operatorname{Circ}(I)$,} qui peut s'interpréter comme \deleted{l'application} \inserted{l'ensemble} \deleted{\unsure{de l'appl.}} \inserted{des} faces d'un cube\ldots\ ].

L'application de degré $2$ $(*)$ définit un cube \inserted{comb.}, dont l'ens. des faces est $\operatorname{Circ}(I)$, \inserted{avec l'antipodisme $u\mapsto u^{-1}$,} l'ens des \deleted{sommets} \deleted{\unsure{arêtes}} arêtes \inserted{(\unsure{resp$^{\mathrm t}$|compr$^{\mathrm t}$|compl.} des sommets)} est l'ens des $12$ sections partielles \deleted{\unsure{dans}} sur des parties de card. $2$ de $J$ (resp. des \deleted{\unsure{$8$|$\mathcal P$} \unsure{les}} $8$ sections de $\operatorname{Circ}(I)$ sur $J$), l'incidence entre facettes \deleted{\unsure{étant}} de dim $2,1,0$ étant le prolongement des \inserted{sections partielles.} \deleted{\unsure{À tout}} Il y aurait

% slanted
\margin{Bien sûr, $G$ opère sur le cube par automorphismes}

% archive sheet number: 22

\scanpage{24}

lieu d'expliciter

a) la correspondance entre \inserted{les $8$} sommets du cube, et les \inserted{$8$} éléments d'ordre $3$ de $G$ (``rotations-sommets'')

b) la \deleted{\unsure{relation|éléments}} \inserted{corr.} entre les $6$ paires d'arêtes antipodiques du cube, et les $6$ \deleted{\unsure{réfle|réfl}} réflexions de $G\unsure{=|\simeq}\unsure{\mathfrak S|G}_I$

c) la corr. entre les $6$ \deleted{éléments de $\operatorname{Circ}(I)$, et} \inserted{faces du cube, et} les \deleted{$6$ \unsure{autres} éléments de $G$} \inserted{permutations circulaires} (él. d'ordre $4$ de $G\simeq\operatorname{Aut}^+(\text{Cube})$) correspondants,

et de décrire également une \uline{orientation canonique} du cube, ou ce qui revient -- $=$, \unsure{une|un} \deleted{trivialisat$^{\mathrm n}$} iso.
\[
\underset{\substack{\Vert\,\mathrm{déf}\\\underline\omega_J\,(\simeq\underline\omega_I)}}{\underline\omega_{\deleted{\unsure{I|J|Q}}}}
\simeq\bigwedge_{i\in J}
\underset{\substack{\text{ens des}\\\text{deux }\deleted{\unsure{\text{élémt}}}\ \inserted{\text{autom.}}\text{ d'ordre }4\\
\text{de la structure carrée}\\Q_i\text{ définie par }i\in J\\
\text{i.e. ens des deux orientations}\\\text{de }Q_i}}{\underbrace{\underline\omega_i}}
\]

Comme on a
\[
\underline\omega_i\simeq
\underset{\substack{\deleted{\unsure{\text{ou}}}\,\text{ens.}\\\text{des deux}\\
\text{codiagonales}\\\deleted{\unsure{\text{du carré}}}\ \inserted{\text{de }Q_i}}}{\underbrace{\underline\omega}}
\wedge
\underset{\substack{\text{ens des diag.}\\\text{pour le carré }Q_i,\ \text{i.e. }D_i=I/\underset{\substack{\text{translation}\\\text{indicée par }i}}{\underbrace{\unsure{\underline a|\underline\sigma}_i}}}}{\underbrace{D_i}}
\]
la formule s'écrit aussi
\[
\underset{\substack{\deleted{\unsure{\text{ens}}}\\\text{torseur}\\\text{sous }\mathbb F_2^J}}{\underbrace{\bigwedge_{i\in J}D_i}}\simeq\mathord{1\mkern-4mu\rule{.055em}{.66em}}
\]

Or on a \deleted{la suite exacte} l'inclusion
\diagram{\begin{tikzcd}[ampersand replacement=\&,row sep=0pt]
I\arrow[r,hook]\&\displaystyle\prod_{i\in J}D_i\\
\substack{\text{torseur}\\\text{sous }V_J(\mathbb F_2)}\&\substack{\subset\ \text{torseur sous}\\\mathbb F_2^J}
\end{tikzcd}}

% archive sheet number: 23

\scanpage{25}
\mspage{12}

qui fait une restriction du groupe structural à $V_J(\mathbb F_2)$, d'où $\bigwedge_{i\in J}D_i\simeq\mathord{1\mkern-4mu\rule{.055em}{.66em}}$, OK.

Il y aurait lieu aussi d'interpréter les $10$ classes d'isomorphie d'espaces homogènes $\ne(1)$ sous $G\simeq\operatorname{Aut}(\text{Cube})$ (correspondant à des \deleted{\unsure{deux|des}} sous-groupe d'isotropie $\ne G$ \unsure{\}|)}, définis à conjugaison près) en termes d'\deleted{objets} éléments de structure associés au cube (tels sommets, arêtes, faces, ou \deleted{tels} objets mod antipodisme -- i.e. sur le cube gauche -- ou \deleted{\unsure{couples de points|repères de points}} \inserted{repères} directs -- ce qui en fait déjà sept \inserted{\unsure{non iso|en iso.}}, de cardinaux $8,12,6,4,6,3,24$. Il y en a un de cardinal $2$, quotient de l'ens des sommets (correspondant à l'ens. $\underline\omega_I$), un autre d'ordre $12$ correspondant aux \deleted{\unsure{fa}} diagonales des faces \inserted{(ou codiagonales, si on préfère)} (les stabilisateurs \deleted{\unsure{de ces él. d'ordre}} \inserted{correspondent} aux él. d'ordre $2$ de $G$ qui \deleted{\unsure{ne}} sont \deleted{\unsure{pas}} dans $V$), enfin un troisième qui est formé des paires d'éléments distincts de $J$ (i.e. repères de $J$), $J$ étant interprété comme l'ens des ``directions de \deleted{\unsure{plans}} \inserted{faces}'' (\deleted{\unsure{ou ens}} (ens de faces/antipodisme) ou des directions d'arêtes. Le compte y est !

% slanted
\margin{position relative de deux sommets à dist. comb. $1$, ou encore de deux arêtes : dist. comb. $1$, \deleted{ou à dist. comb $2$,}}

\deleted{\unsure{5|En}}

% archive sheet number: 24

\scanpage{26}

5. Modules libres de rang $2$ sur $\Lambda=\mathbb Z/4\mathbb Z$.

Soit $M$ un tel module, et soit
\[
M_0=M/2M=M\otimes_\Lambda\mathbb F_2 \tag*{(1)}
\]
qui est un vectoriel de rang $2$ sur $\mathbb F_2$. Soit $J=M_0^*$, on utilise la notation indicielle pour les éléments de $M_0^*$, notés $x_i$ $(i\in J)$. \deleted{Pour tout $i\in J$, soit}

On a une suite exacte de \deleted{groupes comm.} \inserted{$\Lambda$-modules}.
\diagram{\begin{tikzcd}[ampersand replacement=\&,column sep=small,row sep=small]
(2)\&0\arrow[r,no head]\&M_0\arrow[r,no head]\&M\arrow[r,no head,"\mathrm{can}"]\&M_0\arrow[r]\&0\\
\&\&M\arrow[u,dotted,"\mathrm{can}"]\arrow[ur,dotted,"2"]\&\&\&
\end{tikzcd}}
donc $M$ apparaît comme extension de $M_0$ par $M_0$ (en tant que $\Lambda$-module). Pour tt $x\in M_0$, soit $M_x$ son image inverse dans $M$, qui est un torseur sous $M_0$ (i.e. un ens. à quatre éléments \inserted{$M_x$}, muni d'un iso. $M_0\simeq V_{M_x}$).

On posera, pour $i\in J$
\[
\left\{\begin{aligned}
M_i&=M_{x_i}=\deleted{\unsure{\text{image inverse de }\{x_i\}}}\ \inserted{\text{fibre de }M\text{ sur }}x_i\in M_0,\ M_0\text{-torseur}\\
M_J&=\prod_{i\in J}M_i,\quad\text{un torseur sous }\deleted{\unsure{\mathbb F^J|\mathbb F}}\ M_0^J.
\end{aligned}\right. \tag*{(3)}
\]
Comme on a
\[
\sum_{i\in J}x_i=0 \tag*{(4)}
\]
la structure de groupe sur $M$ donne un iso. can.
\[
\bigwedge_{i\in J}M_i=0 \tag*{(6)}
\]
en d'autres termes, introduisant
\[
V_J(M_0)=\{(\deleted{\unsure{x}}\unsure{\xi|\zeta}_i)_{i\in J}\in M_0^J\mid\sum\unsure{\zeta|\xi}_i=0\} \tag*{(7)}
\]
s'insérant dans la suite exacte
\diagram{\begin{tikzcd}[ampersand replacement=\&,column sep=small]
(8)\&0\arrow[r,no head]\&V_J(M_0)\arrow[r,no head]\&M_0^J\arrow[r,no head,"\text{somme}"']\&M_0\arrow[r,no head]\&0,
\end{tikzcd}}
le torseur $M_J$ sous $M_0^J$ provient canoniquement d'un torseur sous $V_J(M_0)$, que je note $P_M$,

% archive sheet number: 25

\scanpage{27}
\mspage{13}

De façon précise, \inserted{la donnée} des \inserted{$M_0$-torseurs} $M_i$ et \inserted{de} l'isom can. \unsure{(6)} revient à la donnée du torseur $P_M$ sous $V_J(M_0)$, avec
\[
M_i\simeq P_M\wedge_{M_0^J}(M_0,\operatorname{pr}_i)\qquad\text{pour }i\in J \tag*{(9)}
\]
où $\operatorname{pr}_i:M_0^J\mathrel{\text{---}}M_0$ la projection d'indice $i\in J$.

Notons qu'on a des iso. canoniques (valables pour tt vectoriel de rang $2$ sur $\mathbb F_2$)
\[
\begin{aligned}
\operatorname{End}(M_0)&\simeq V_J(M_0)\xrightarrow{\sim}M_0\otimes M_0\simeq M_0\oplus\mathbb F_2^{\underline\omega_{M_0}}\\
\underline\omega_{M_0}&\simeq\underline\omega(\underset{J}{\underbrace{M_0^*}})\quad\text{ens des orientations de }J.
\end{aligned}\tag*{(10)}
\]

% slanted
\margin{\uline{NB} Pour tt vectoriel $W$ sur $\mathbb F_2$ on a (11) $\operatorname{Hom}_{\mathbb F_2}(M_0,W)\simeq V_J(\unsure{W|V})$ Car $M_0\simeq\mathbb F_2^J/\operatorname{diag}$ \deleted{est} $\longrightarrow$}

\[
\left\{\begin{aligned}
\operatorname{End}(M_0)&\xrightarrow{\sim}V_J(M_0)\\
u&\longmapsto(u(x_i))_{i\in J}
\end{aligned}\right. \tag*{(12)}
\]
\[
\operatorname{End}(M_0)\simeq M_0^\vee\otimes M_0\simeq M_0\otimes M_0
\qquad\substack{\text{car }M_0^\vee\simeq M_0\\\text{(valable pour tt}\\\text{vect. de rang }2\text{ sur }\mathbb F_2)}
\]

\diagram{\begin{tikzcd}[ampersand replacement=\&,column sep=small,row sep=0pt]
(13)\& M_0\arrow[r,hook]\&\operatorname{End}(M_0)\\
\&\rotatebox{90}{$\subset$}\&\rotatebox{90}{$\subset$}\\
\&M_0^*\arrow[r,leftrightarrow]\&\operatorname{Aut}(M_0)\simeq\unsure{\mathfrak S|G}_J\\
\&x_i\arrow[r,mapsto]\&\substack{\deleted{\text{transposition}}\ \inserted{\sigma_i}\\\text{réflexion de }M_0\text{ qui fixe }x_i\\
\text{(donc }\deleted{\unsure{\text{qui échange}}}\ \inserted{\text{transposition}}\text{ de }x_j\text{ et }x_k)}
\end{tikzcd}}

[\uline{NB} Comme $\sum_{i\in J}\sigma_i=0$, cela correspond à \unsure{hom}
\[
M_0\longrightarrow\operatorname{End}(M_0)\ ]
\]

\diagram{\begin{tikzcd}[ampersand replacement=\&,column sep=small,row sep=0pt]
(14)\&\mathbb F_2^{\underline\omega_{M_0}}\arrow[r,hook]\&\operatorname{End}(M_0)\\
\&\rotatebox{90}{$\subset$}\&\rotatebox{90}{$\subset$}\\
\&\underline\omega_{M_0}\arrow[r,hook,"\mathrm{can}"]\&\operatorname{Aut}(M_0)\simeq\mathfrak S_J\\
\&\Vert\&\\
\&\underline\omega_J\simeq\mathfrak S_J^+\smallsetminus\{\operatorname{id}_J\}\&
\end{tikzcd}}

\uline{NB} \deleted{\unsure{Comme}} \inserted{On a} \inserted{si $\rho\in\underline\omega$, $\rho'\overset{\mathrm{déf}}=\rho^{-1}$}
\[
\rho+\rho'\deleted{\unsure{=|+}}\,+\unsure{\deleted{1}|1}\operatorname{id}=0
\]
i.e. $\rho+\rho'=\operatorname{id}_{M_0}$

i.e. l'hom $\mathbb F_{\unsure{2|1}}^{\underline\omega}\mathrel{\unsure{\lhook\joinrel\relbar|\hookrightarrow}}\operatorname{End}(M_0)$ applique $1$ sur $1$.

Si $i,j\in J$, $i\ne j$, alors $\sigma_i,\sigma_j,\rho,\rho'\in\operatorname{End}(M_0)$ sont lin. indépendants sur $\mathbb F_2$ donc forment une base, \deleted{et} alors que \unsure{$x_i,x_j$|$\sigma_i,\sigma_j$|\deleted{$x_i,x_j$}} en forment une de $M_0$, et $\rho,\rho'$ de $\mathbb F_2^{\underline\omega}$ \ldots

% archive sheet number: 26

